\section{Gabarito e resoluções comentadas — Direito Constitucional}

\begin{solucao}{05.01}\textbf{CERTO.} Essa é a classificação doutrinária tradicional da CF/1988.\end{solucao}
\begin{solucao}{05.02}\textbf{ERRADO.} Soberania pertence à República Federativa do Brasil; os entes federados possuem autonomia.\end{solucao}
\begin{solucao}{05.03}\textbf{CERTO.} Direitos podem colidir e sofrer restrições constitucionalmente justificadas; não existe absolutismo geral.\end{solucao}
\begin{solucao}{05.04}\textbf{ERRADO.} Nacionalidade é vínculo jurídico-político com o Estado; cidadania envolve exercício de direitos políticos.\end{solucao}
\begin{solucao}{05.05}\textbf{CERTO.} São entes autônomos, cada qual com competências constitucionais.\end{solucao}
\begin{solucao}{05.06}\textbf{ERRADO.} Tribunais de Contas não integram o Judiciário.\end{solucao}
\begin{solucao}{05.07}\textbf{CERTO.} Há reserva de matéria e quórum de maioria absoluta para lei complementar.\end{solucao}
\begin{solucao}{05.08}\textbf{ERRADO.} Não existe hierarquia geral entre LC e LO; a Constituição distribui matérias e procedimentos.\end{solucao}
\begin{solucao}{05.09}\textbf{CERTO.} O CNJ exerce controle administrativo/financeiro e disciplinar nos limites constitucionais, não jurisdição recursal típica.\end{solucao}
\begin{solucao}{05.10}\textbf{CERTO.} PPA, LDO e LOA integram o ciclo constitucional orçamentário.\end{solucao}
\begin{solucao}{05.11}\textbf{B.} Dignidade da pessoa humana é fundamento do art. 1º. Erradicação da pobreza e desenvolvimento são objetivos; defesa da paz e cooperação são princípios internacionais.\end{solucao}
\begin{solucao}{05.12}\textbf{B.} Restrições precisam de fundamento e exame de proporcionalidade, sem destruir o núcleo protegido.\end{solucao}
\begin{solucao}{05.13}\textbf{CERTO.} Habeas data protege acesso/retificação de dados pessoais nas hipóteses constitucionais e legais.\end{solucao}
\begin{solucao}{05.14}\textbf{C.} Mandado de injunção enfrenta omissão normativa que inviabiliza direito ou prerrogativa constitucional.\end{solucao}
\begin{solucao}{05.15}\textbf{B.} A competência concorrente é legislativa e repartida entre União, Estados e DF segundo o art. 24.\end{solucao}
\begin{solucao}{05.16}\textbf{CERTO.} Intervenção é exceção constitucionalmente tipificada à autonomia.\end{solucao}
\begin{solucao}{05.17}\textbf{B.} Concurso é a regra para cargo/emprego, ressalvadas hipóteses constitucionais.\end{solucao}
\begin{solucao}{05.18}\textbf{B.} A regra é vedação, com exceções expressas e compatibilidade de horários.\end{solucao}
\begin{solucao}{05.19}\textbf{CERTO.} A Constituição prevê hipóteses de perda do cargo do servidor estável.\end{solucao}
\begin{solucao}{05.20}\textbf{B.} Separação dos Poderes integra o núcleo protegido contra proposta tendente à abolição.\end{solucao}
\begin{solucao}{05.21}\textbf{B.} No plano federal, o Congresso exerce controle externo com auxílio do TCU.\end{solucao}
\begin{solucao}{05.22}\textbf{CERTO.} Essa é distinção constitucional central entre parecer prévio das contas presidenciais e julgamento das contas dos responsáveis.\end{solucao}
\begin{solucao}{05.23}\textbf{A.} MP exige relevância e urgência e possui limites materiais/procedimentais.\end{solucao}
\begin{solucao}{05.24}\textbf{B.} O CNJ controla administração, finanças e deveres funcionais do Judiciário, nos termos da CF.\end{solucao}
\begin{solucao}{05.25}\textbf{CERTO.} A ação popular é remédio constitucional de titularidade do cidadão para hipóteses constitucionalmente protegidas.\end{solucao}
\begin{solucao}{05.26}\textbf{C.} I e III são corretas; aplicação imediata não significa inexistência de qualquer necessidade de conformação normativa.\end{solucao}
\begin{solucao}{05.27}\textbf{CERTO.} Igualdade admite diferenciação juridicamente justificada e proporcional.\end{solucao}
\begin{solucao}{05.28}\textbf{A.} Assuntos de interesse local são, em regra, competência municipal.\end{solucao}
\begin{solucao}{05.29}\textbf{B.} Estados detêm competências remanescentes não vedadas pela Constituição.\end{solucao}
\begin{solucao}{05.30}\textbf{CERTO.} Art. 23 trata competências materiais comuns; art. 24, legislativas concorrentes.\end{solucao}
\begin{solucao}{05.31}\textbf{B.} Estabilidade pressupõe cargo efetivo e requisitos constitucionais; não alcança automaticamente comissionados.\end{solucao}
\begin{solucao}{05.32}\textbf{B.} Lei complementar exige maioria absoluta.\end{solucao}
\begin{solucao}{05.33}\textbf{CERTO.} Vício de iniciativa/reserva formal pode gerar inconstitucionalidade formal.\end{solucao}
\begin{solucao}{05.34}\textbf{B.} O TC auxilia o Legislativo, mas possui posição institucional e competências próprias, sem subordinação hierárquica.\end{solucao}
\begin{solucao}{05.35}\textbf{B.} É a distinção paradigmática estabelecida na Constituição federal.\end{solucao}
\begin{solucao}{05.36}\textbf{CERTO.} O art. 70 inclui legalidade, legitimidade e economicidade, entre outros parâmetros.\end{solucao}
\begin{solucao}{05.37}\textbf{B.} Controle difuso ocorre incidentalmente em casos concretos por órgãos jurisdicionais competentes.\end{solucao}
\begin{solucao}{05.38}\textbf{B.} ADI perante o STF é instrumento do controle concentrado federal.\end{solucao}
\begin{solucao}{05.39}\textbf{CERTO.} Precedentes, modulação e mecanismos constitucionais/legais tornam os efeitos mais complexos que a fórmula escolar absoluta.\end{solucao}
\begin{solucao}{05.40}\textbf{B.} A LOA estima receitas e fixa despesas.\end{solucao}
\begin{solucao}{05.41}\textbf{CERTO.} O auxílio ao Legislativo não cria relação hierárquica sobre decisões constitucionalmente atribuídas ao Tribunal.\end{solucao}
\begin{solucao}{05.42}\textbf{B.} Economicidade integra expressamente o alcance do controle constitucional.\end{solucao}
\begin{solucao}{05.43}\textbf{B.} A igualdade constitucional exige examinar justificativa e proporcionalidade da diferenciação.\end{solucao}
\begin{solucao}{05.44}\textbf{CERTO.} Contraditório efetivo envolve possibilidade real de manifestação e influência, quando devido.\end{solucao}
\begin{solucao}{05.45}\textbf{A.} O primeiro filtro é a repartição constitucional de competências.\end{solucao}
\begin{solucao}{05.46}\textbf{CERTO.} Entes federados são autônomos, não soberanos.\end{solucao}
\begin{solucao}{05.47}\textbf{B.} Usar espécie inadequada em matéria reservada pode configurar vício formal.\end{solucao}
\begin{solucao}{05.48}\textbf{B.} Tribunais de Contas possuem competências decisórias e sancionatórias nos limites constitucionais/legais.\end{solucao}
\begin{solucao}{05.49}\textbf{CERTO.} A cláusula pétrea impede proposta tendente a abolir o núcleo, não qualquer alteração textual relacionada ao tema.\end{solucao}
\begin{solucao}{05.50}\textbf{A.} A constitucionalidade depende da competência legislativa e de eventual mecanismo constitucional de delegação.\end{solucao}
